% !TeX root = ../papers/nc-p3.tex
\section{有限点枝の反対階数八と両対称性}
\label{sec:rank-eight-finite-symmetry}
\begin{theorem}\label{thm:rank-eight-finite-symmetry}
標準有限Frobenius窓の四辺が有限、反対階数が8で、長さ20の点スキームに長さ1の局所成分が一つあるとする。
中支持3の点があれば全次数で正則である。それ以外の窓は全て独立な基底変更で両対称反転窓になる。
従ってこの範囲では全次数の正則性と添字保持同型類は有限に決まる。両対称の場合は最小周期が1または2で、
純関係と大きな射影ファイバーは既に排除されているので、残る正則性の検査は二変数Hankel商の不在だけである。
\end{theorem}
\begin{proof}[二つの交換子空間]
系\ref{cor:point-commutator-rank}の印付き表示を使う。中支持3の場合は定理\ref{thm:finite-cubic-separators}による。
それ以外では $G,F$ は同じ既約な跡零三平面 $\mathcal G\subset\operatorname{End}(V)$ の二基底、
$Q,P$ は同様の $\mathcal Q\subset\operatorname{End}(U)$ の二基底である。
既約とは共通不変部分空間が非自明にはないという意味である。$F_c=G_{Tc}$、$P_d=Q_{Jd}$ と書き、
交換子式を
\[
 [G_a,G_{Tc}]_{bd}=-[Q_b,Q_{Jd}]_{ac}
\]
とする。$T,J$ は可逆である。この式の $ac\mid bd$ 階数 $r$ は両三平面の交換子の線型包の次元に等しい。
$r\le1$ なら $\wedge^2\mathcal G$ の核の二次元部分が共通因子を持ち、
その非零因子が $\mathcal G$ 全体と交換する。Schurの補題と跡零性に反する。
従って $r=2,3$ である。さらに各交換子空間は、三行三列の交代行列空間の片側を可逆変換したものに入る。
実際、右辺は $b,Jd$ について交代し、左辺にも同じ議論が適用できる。
\end{proof}

\begin{proof}[交換子の次元が二の場合]
$\mathcal G$ 内の二平面の基底 $X,Y$ が可換であり、第三の基底を $Z$ とする。
その可換平面に最小多項式が三次の $X$ を取れる。実際、二固有値の半単純元しか三次に届かない
可能性を考えると、その二次元固有空間上で別の可換元がスカラーでなければ、一般線型結合の
最小多項式は三次になる。スカラーなら跡零性から二元が従属となる。
よって三次の元がなければ全元が平方零となる（単一固有値は跡零なので0）。
この二平面は全て階数1以下で、共通像直線または共通核平面を持つ。
前者なら $[X,Z],[Y,Z]$ の像はその直線と $Z$ での像の和に入り、後者なら共通核とその $Z$ 逆像の交わりを消す。
これは独立な二つの交代三行列が全体で三次元を張る像と自明な共通核を持つことに反する。
この二性質は上の片側可逆変換でも保たれる。

巡回行列 $X$ に可換な $Y$ はその多項式なので $(X,Z)$ は既約である。
$\det[X,Z]=0$ は、ある非零対称行列 $H$ が $HX=X^tH,\ HZ=Z^tH$ を満たすことと同値である。
実際、対称行列六次元から交代行列二組の六次元へのこの線型写像の行列式は符号を除き $\det[X,Z]$ である。
一般の $X$ を対角化すると、三つの非対角未知数の係数は固有値差で、対角未知数の残りの
行列式は $Z_{12}Z_{23}Z_{31}-Z_{13}Z_{32}Z_{21}$ の符号倍である。
同時共役で行列式の変換倍率は1となる（対称二乗の行列式の指標は四乗、交代二乗二組も四乗）。
多項式恒等式なので全域で成り立つ。ここでは交換子空間が交代行列の片側変換に入るため行列式は零である。
$H$ の核は不変なので $H$ は可逆で、$Y$ も同じ形式に関して自己随伴である。
$\mathcal Q$ にも適用し、二空間の基底を独立に変えれば四組 $G,F,Q,P$ はいずれも対称行列になる。

このとき交換子式の右辺は行列として $a,c$ の交換で符号を変える。従って全ての $a$ で $[G_a,G_{Ta}]=0$。
可換二平面の外の $G_a$ と交換する $\mathcal G$ の元は $G_a$ の倍しかない
（交換子写像の核はこの平面の外積の直線である）。従って $T$ はスカラーで、$J$ も同様である。
\end{proof}

\begin{proof}[交換子の次元が三の場合]
交換子式を $a,c$ で跡を取れば $\sum_a[G_a,G_{Ta}]=0$。三つの交換子の独立性により $T$ は対称で、
$J$ も対称である。可逆対称形式をそれぞれ合同変換で単位行列に取り、$F=G,P=Q$ とする。
このとき $[G_a,G_c]$ と $[Q_b,Q_d]$ は交代行列である。

三基底 $G_i=S_i+K_i$ を跡零対称部分と交代部分に分けると $[K_i,S_j]=[K_j,S_i]$ である。
$K_i=[v_i]_\times$ は三次元外積作用とする。$v_i$ の張る次元が1なら二つの独立な対称部分が
同じ非零外積行列と交換することになる。しかしその行列は巡回行列で（等方ベクトルでも平方が非零の冪零行列）、
跡零対称な中心化空間は一次元である。次元2なら第三の非零対称部分が独立な二つの外積行列と交換する。
その平方 $vv^t-(v^tv)I$ とも交換するため、二平面内の各 $v$ はその固有ベクトルである。
従ってこの対称部分はスカラーを除くと $cww^t$、$w$ は平面に垂直となる。
$c\ne0$ なら交換関係から全ての $v$ で $v\times w=0$ となり矛盾する。
次元3なら三基底を線型結合して $v_i=e_i$ とする。上の三つの交換子式は $S_i$ の十五成分に非退化である。
具体的に各対称跡零行列に基底
$E_{11}-E_{33},E_{22}-E_{33},E_{12}+E_{21},E_{13}+E_{31},E_{23}+E_{32}$ を取り、
対 $23,31,12$ の順に差の成分 $11,22,12,13,23$ を読むと行列式は $-3456$ である。
従って $\mathcal G$ は全て対称か全て交代であり、$\mathcal Q$ も同様である。

混合型は反対階数条件に反する。$G$ を交代側として $G_i=[Le_i]_\times$ と書き、
定理\ref{thm:marked-opposite-eight}と同じ向きで交換子行列 $D_G$ と
$P_G(Z)_{ac}=\operatorname{tr}(ZG_aG_c)$ を取る。対称跡零 $Z$ に対して
\[
 D_G P_G(Z)D_G^t=-\det(D_G) Z
\]
である。$D_G=-\operatorname{cof}(L)$ と $P_G(Z)=L^t ZL$ は外積の積の式による。
一方、対称な $Q$ では補題\ref{lem:ternary-net-commutator}により、
$I,Q_1,Q_2,Q_3$ に直交する対称行列 $Y$ に対して同じ式の符号が正である。
$D_Q=-D_G^t$ なので残余対合は
\[
 \operatorname{tr}(Y P_G(Z))+\operatorname{tr}(Z P_Q(Y))
 =-2\det(D_G)\operatorname{tr}(Z D_G^{-t}Y D_G^{-1}).
\]
$Y$ の空間は二次元以上であり、その可逆合同変換の像が全てスカラーになることはないから、
右辺は恒等零ではない。ところが系\ref{cor:point-commutator-rank}の $\bar w=0$ はまさにこれを消す。
逆の混合型も位置を交換すればよい。

双方が交代なら $Q_b=[Ne_b]_\times$ と書いて
$\operatorname{cof}(N)=-\operatorname{cof}(L)^t$ より $N$ は $L^t$ の非零倍となる。
最初の辺に点対 $((r,x),(s,y))$ を入れた六式は
$r y+G_x^t y=0,\ s x+Q_y^t x=0$ である。実際、これらは反対側の星印一つを含む六成分であり、
印付き表示の可逆な六次小行列により独立である。$x,y$ をそれぞれ非零等方ベクトルとし、
$y^tLx=0$ を課す。両等方二次曲線の積上のこの零点は次元1以上である。
$(Lx)\times y$ は $y$ の倍であり、$(Ny)\times x$ も $x$ の倍なので
$r,s$ が取れて辺の実際の点対となる。これは有限性に反する。
\end{proof}

\begin{proof}[両対称表示]
残る全ての場合に $G,Q$ が対称で $F=hG,P=kQ$、$h,k\ne0$ である。
位置0,1の星印の係数をそれぞれ $1/h,1/k$ 倍すれば、星印を持つ項は各対向位置の交換で不変となる。
純成分 $h(G_aG_c)_{bd}+k(Q_bQ_d)_{ac}$ も交換子式で各交換に不変である。
よって四テンソルは両対称であり、定理\ref{thm:bisymmetric-two-length-test}により結論を得る。
同型・周期は定理\ref{thm:finite-all-window-keys}で決まる。
\path{computations/verify_rank_eight_finite_symmetry.py}は十五成分の行列式、六未知数の対称化方程式、
および二種の符号のうち交代側を検算する。
\end{proof}
